% main.tex — Informe técnico del proyecto final de Gestión de Datos (UAX, 2025/26)
% Autor: Alejandro Barreche
%
% Compilación:
%   latexmk -pdf -interaction=nonstopmode main.tex
% o, alternativamente:
%   pdflatex -interaction=nonstopmode main.tex (dos pasadas)
%
% No hay bibliografía: el informe se sostiene sobre el código y los datos
% del propio repositorio. No requiere biber.

\documentclass[11pt, a4paper, oneside]{report}
\usepackage{projectreport}

% --- Metadatos del informe ---------------------------------------------------
\renewcommand{\name}{Alejandro Barreche}
\renewcommand{\course}{Grado en Ingeniería Matemática}
\renewcommand{\projecttitle}{%
  SalesHealth: construcción de un Data Warehouse y segmentación de clientes
  mediante CLTV y PCA+KMeans%
}
\renewcommand{\submissiondate}{Mayo 2026}

\begin{document}

% ------------------------------------------------------------------ Portada --
\makecoverpage

% --------------------------------------------------------- Resumen / abstract
\chapter*{Resumen}
\addcontentsline{toc}{chapter}{Resumen}
\thispagestyle{plain}

Este informe documenta el proyecto \textbf{SalesHealth}, entrega final de la
asignatura Gestión de Datos del curso 2025/26 en la Universidad Alfonso X El
Sabio. Sobre un dataset relacional sintético del sector salud y bienestar
($5\,750$ clientes, $20\,000$ tickets de venta, $42\,555$ líneas, $2\,330$
devoluciones, ventana 2020-01--2026-01) se construye un \emph{Data Warehouse}
con esquema en estrella sobre PostgreSQL, se calcula el \emph{Customer
Lifetime Value} (CLTV) por cliente según la fórmula del enunciado
$\mathrm{CLTV} = \text{Ingresos}_t \times \text{Margen}_t \times
\text{Frecuencia}_t \times R_t$, y se segmenta la base mediante PCA y KMeans
sobre ocho métricas derivadas. La cadena ETL queda automatizada por seis
notebooks numerados y por SQL idempotente; los resultados se consumen desde
un dashboard estático HTML+Plotly servido por Flask. El proyecto incluye una
batería de pruebas \texttt{pytest} para la integridad referencial del DWH y
para la consistencia de la fórmula de CLTV.


% --------------------------------------------------------------- Índices ----
\tableofcontents

% --------------------------------------------------------------- Capítulos --
\section{Introducción}
\label{sec:intro}

El proyecto parte de un dataset sintético del sector salud y bienestar
distribuido como punto de partida para la entrega final. Reproduce un
escenario retail clásico ($50$ productos, una red de tiendas, catálogo de
promociones, devoluciones con motivo) sobre una ventana de $72$ meses
($2020$--$2025$, ambos inclusive). Los datos crudos están cargados en el
esquema \texttt{public} de PostgreSQL~16 en Docker (\texttt{localhost:5433},
BD \texttt{saleshealth}). La tabla~\ref{tab:cifras-clave} reúne las cifras
verificadas contra el DWH y los \emph{parquet} de analítica.

\begin{table}[H]
\centering
\caption{Cifras clave verificadas (DWH + \texttt{data/processed/}).}
\label{tab:cifras-clave}
\small
\begin{tabular}{lr@{\hspace{1.4em}}lr}
\toprule
Productos          & $50$            & Líneas de venta   & $42\,555$       \\
Clientes           & $5\,750$        & Líneas devueltas  & $2\,330$        \\
Tickets de venta   & $20\,000$       & Ingresos brutos   & $9{,}68$\,M\texteuro \\
Ingresos netos     & $9{,}40$\,M\texteuro & Ventana          & $72$ meses    \\
\bottomrule
\end{tabular}
\end{table}

\noindent\textbf{Objetivos.} (i) Diseñar un modelo Entidad--Relación sobre
\texttt{public} y un modelo dimensional en estrella sobre \texttt{dwh}.
(ii) Implementar un ETL reproducible \texttt{public$\to$staging$\to$dwh}.
(iii) Calcular el CLTV por cliente,
$\mathrm{CLTV}=\text{Ingresos}_t \times \text{Margen}_t \times
\text{Frecuencia}_t \times R_t$, junto con \emph{Recencia}, \emph{AOV} y
\emph{tasa de devolución}. (iv) Segmentar la base mediante \emph{PCA} +
\emph{KMeans}. (v) Documentar todo en un informe y un dashboard.

\noindent\textbf{Stack.} Python~3.11 en entorno \texttt{conda} \texttt{UAX};
PostgreSQL~16; \texttt{pandas}, \texttt{SQLAlchemy}~2.0, \texttt{scikit-learn};
notebooks Jupyter como orquestador; Flask para el dashboard; \texttt{pytest}.

\section{Modelo de datos}
\label{sec:modelo}

El esquema \texttt{public} agrupa $17$ tablas crudas en cinco bloques:
\emph{producto} (\texttt{product},\allowbreak{} \texttt{central\_product},\allowbreak{}
\texttt{category}, \texttt{brand}, \texttt{offer},\allowbreak{}
\texttt{product\_offer}); \emph{almacén} (\texttt{warehouse},\allowbreak{}
\texttt{warehouse\_location},\allowbreak{} \texttt{central\_inventory},\allowbreak{}
\texttt{inventory}); \emph{tienda y geografía} (\texttt{store},
\texttt{city\_zone}); \emph{cliente y ventas} (\texttt{customer},
\texttt{sale}, \texttt{sale\_item}); y \emph{devoluciones}
(\texttt{return\_item}, \texttt{return\_reason}). Las relaciones clave
son \texttt{sale}\,1--N\,\texttt{sale\_item}, \texttt{sale}\,N--1
\texttt{customer} y \texttt{store}, \texttt{sale\_item}\,N--1
\texttt{product}, y \texttt{sale\_item}\,1--N\,\texttt{return\_item}
(figura~\ref{fig:er}).

\begin{figure}[H]
  \centering
  \includegraphics[width=0.92\linewidth]{er_diagram.png}
  \caption{Esquema \texttt{public}: $17$ tablas agrupadas en cinco bloques
  funcionales.}
  \label{fig:er}
\end{figure}

El esquema \texttt{dwh} sigue un patrón clásico en estrella con seis
dimensiones (\texttt{dim\_date}, \texttt{dim\_customer},
\texttt{dim\_product}, \texttt{dim\_store}, \texttt{dim\_offer},
\texttt{dim\_return\_reason}) y dos hechos: \texttt{fact\_sales}
(grano $=$ línea de venta, $42\,555$~filas, métricas \texttt{quantity},
\texttt{unit\_price}, \texttt{unit\_cost}, \texttt{subtotal},
\texttt{cost}, \texttt{margin}) y \texttt{fact\_returns}
(grano $=$ línea devuelta, $2\,330$~filas, métricas \texttt{quantity} y
\texttt{returned\_value}). Véase la figura~\ref{fig:dim}.

\begin{figure}[H]
  \centering
  \includegraphics[width=0.92\linewidth]{dim_diagram.png}
  \caption{Esquema \texttt{dwh}: seis dimensiones más dos hechos. La línea
  discontinua representa el \emph{drill-through} de
  \texttt{fact\_returns} hacia \texttt{fact\_sales} por
  \texttt{sale\_item\_id}.}
  \label{fig:dim}
\end{figure}

\subsection*{Decisiones de diseño}

Todas las dimensiones tienen una \emph{surrogate key} \texttt{*\_sk}
(\texttt{serial}) y conservan la clave de negocio como \texttt{unique},
lo que hace la carga idempotente (\texttt{truncate \dots restart
identity}). \texttt{dim\_product} se construye con \texttt{LEFT JOIN}
entre \texttt{product} ($50$~filas) y \texttt{central\_product}
($49$~filas) — un producto carece de coste y su margen se gestiona con
\texttt{coalesce(unit\_cost,\,0)}. \texttt{fact\_returns} desnormaliza
\texttt{customer\_sk}, \texttt{product\_sk} y \texttt{store\_sk} desde
\texttt{fact\_sales} a través de \texttt{sale\_item\_id}, lo que evita
\emph{re-joins} en analítica de devoluciones. En el área \texttt{staging}
se recomputa \texttt{subtotal} en el $100\%$ de las líneas porque $8$
filas raw eran inconsistentes con $\texttt{quantity}\times
\texttt{unit\_price}$; el test
\texttt{test\_revenue\_matches\_between\_staging\_and\_dwh} certifica
que la corrección no rompe la conservación de ingresos. El DDL completo
está en \texttt{sql/02\_dwh\_schema/}.

\chapter{ETL: del raw al DWH}
\label{ch:etl}

\section{Pipeline y orquestación}

El proyecto adopta una arquitectura ETL clásica en tres saltos
(figura conceptual):

\[
  \underbrace{\texttt{public}}_{\text{raw}}
  \;\longrightarrow\;
  \underbrace{\texttt{staging}}_{\text{cleaned}}
  \;\longrightarrow\;
  \underbrace{\texttt{dwh}}_{\text{star schema}}
\]

La orquestación es deliberadamente simple: seis notebooks Jupyter
numerados (\texttt{notebooks/00\_\dots} a \texttt{notebooks/05\_\dots}) que
se ejecutan en orden y comparten utilidades reutilizables alojadas en el
paquete \texttt{src/}. No hay Airflow, Prefect ni Dagster. La
tabla~\ref{tab:notebooks} resume cada paso.

\begin{table}[H]
\centering
\caption{Notebooks que orquestan el pipeline ETL y de analítica.}
\label{tab:notebooks}
\small
\begin{tabular}{l l p{7.5cm}}
\toprule
\textbf{\#} & \textbf{Notebook} & \textbf{Salida} \\
\midrule
00 & \texttt{00\_exploration.ipynb}      & Perfilado de \texttt{public}.        \\
01 & \texttt{01\_etl\_staging.ipynb}     & Tablas \texttt{staging.*} limpias.   \\
02 & \texttt{02\_etl\_dwh.ipynb}         & Esquema \texttt{dwh} cargado.        \\
03 & \texttt{03\_cltv\_calculation.ipynb} & \texttt{data/processed/cltv.parquet}.\\
04 & \texttt{04\_customer\_metrics.ipynb}&
   \texttt{data/processed/customer\_metrics.parquet}. \\
05 & \texttt{05\_pca\_clustering.ipynb}  &
   \texttt{models/\{pca,kmeans,scaler\}.joblib} y
   \texttt{models/customer\_segments.parquet}. \\
\bottomrule
\end{tabular}
\end{table}

\section{Transformaciones SQL}

Las transformaciones que generan \texttt{dwh} a partir de
\texttt{staging} viven como scripts SQL versionados en el directorio
\texttt{sql/03\_transformations/}. Sus nombres siguen un esquema de
prefijos numéricos para fijar el orden de ejecución:

\begin{enumerate}
  \item \texttt{10\_load\_dim\_date.sql}: genera \texttt{dim\_date}
        cubriendo el rango temporal completo de ventas y devoluciones
        mediante \texttt{generate\_series}.
  \item \texttt{20\_load\_dim\_customer.sql}: inserta clientes
        normalizando el nombre completo
        (\texttt{first\_name} + \texttt{last\_name} + \texttt{last\_name2}).
  \item \texttt{30\_load\_dim\_product.sql}: combina \texttt{product} con
        \texttt{central\_product} mediante \texttt{LEFT JOIN}, y resuelve
        \texttt{category} y \texttt{brand} a partir de sus respectivos
        catálogos.
  \item \texttt{40\_load\_dim\_store.sql}: enriquece la tienda con la
        geografía (\texttt{city\_zone}).
  \item \texttt{50\_load\_dim\_misc.sql}: carga \texttt{dim\_offer} y
        \texttt{dim\_return\_reason} en una sola pasada.
  \item \texttt{60\_load\_fact\_sales.sql}: carga el hecho de ventas con
        coste y margen ya calculados.
  \item \texttt{70\_load\_fact\_returns.sql}: carga el hecho de
        devoluciones reusando las FKs de \texttt{fact\_sales} a través
        de \texttt{sale\_item\_id}.
\end{enumerate}

El cuadro~\ref{lst:dim_date} muestra el script más representativo: la
construcción de la dimensión calendario. Genera todas las fechas entre la
primera venta y la última devolución, infiere el resto de atributos
(\texttt{year}, \texttt{quarter}, día de la semana, indicador de fin de
semana, etc.) y los inserta en \texttt{dim\_date}.

\begin{lstlisting}[language=SQLpostgres,
  caption={\texttt{sql/03\_transformations/10\_load\_dim\_date.sql}:
  generación de la dimensión calendario por \texttt{generate\_series}.},
  label={lst:dim_date}]
truncate table dwh.dim_date cascade;

with bounds as (
    select  least(
                (select min(sale_date)::date   from staging.sale),
                (select min(return_date)::date from staging.return_item)
            ) as ts,
            greatest(
                (select max(sale_date)::date   from staging.sale),
                (select max(return_date)::date from staging.return_item)
            ) as te
),
calendar as (
    select generate_series(b.ts, b.te, interval '1 day')::date as d from bounds b
)
insert into dwh.dim_date
select  to_char(d, 'YYYYMMDD')::int        as date_id,
        d                                  as full_date,
        extract(year  from d)::int         as year,
        extract(quarter from d)::int       as quarter,
        extract(month from d)::int         as month,
        to_char(d, 'TMMonth')              as month_name,
        extract(day from d)::int           as day,
        extract(isodow from d)::int        as day_of_week,
        to_char(d, 'TMDay')                as day_name,
        extract(week from d)::int          as week_of_year,
        extract(isodow from d) in (6,7)    as is_weekend
from calendar;
\end{lstlisting}

La carga de \texttt{fact\_sales} ilustra cómo se inyectan las
\emph{surrogate keys} y cómo se calcula \texttt{cost} y \texttt{margin}
sin asumir que todos los productos tengan coste:

\begin{lstlisting}[language=SQLpostgres,
  caption={Carga de \texttt{fact\_sales} reutilizando las dimensiones ya
  pobladas.},
  label={lst:fact_sales_load}]
insert into dwh.fact_sales
       (sale_item_id, sale_id, sale_date_id, customer_sk, product_sk,
        store_sk, offer_sk, quantity, unit_price, unit_cost,
        subtotal, cost, margin)
select  si.sale_item_id,
        s.sale_id,
        to_char(s.sale_date, 'YYYYMMDD')::int                  as sale_date_id,
        dc.customer_sk, dp.product_sk, ds.store_sk, do2.offer_sk,
        si.quantity, si.unit_price, dp.unit_cost,
        si.subtotal,
        coalesce(dp.unit_cost, 0) * si.quantity                as cost,
        si.subtotal - coalesce(dp.unit_cost, 0) * si.quantity  as margin
from staging.sale_item si
join staging.sale          s   on s.sale_id    = si.sale_id
left join dwh.dim_customer dc  on dc.customer_id = s.customer_id
left join dwh.dim_product  dp  on dp.product_id  = si.product_id
left join dwh.dim_store    ds  on ds.store_id    = s.store_id
left join dwh.dim_offer    do2 on do2.offer_id   = si.offer_id;
\end{lstlisting}

\section{Helpers Python}

Los scripts SQL anteriores se ejecutan desde un cuaderno mediante un
helper minimalista alojado en \texttt{src/etl/load.py}: dado un
directorio, ordena los \texttt{.sql} alfabéticamente y los aplica con
\texttt{run\_sql\_file}.

\begin{lstlisting}[language=Python,
  caption={\texttt{src/etl/load.py}: ejecutor de transformaciones SQL.},
  label={lst:run_sql_dir}]
def run_sql_directory(directory: str | Path) -> list[Path]:
    """Ejecuta los .sql de un directorio en orden alfabetico."""
    files = sorted(Path(directory).glob("*.sql"))
    for f in files:
        log.info("Ejecutando %s", f.name)
        run_sql_file(f)
    return files
\end{lstlisting}

\noindent Por debajo, la conexión es un singleton SQLAlchemy expuesto por
\texttt{src/db/engine.py}, configurado desde \texttt{config/settings.py}
(que a su vez lee \texttt{.env}). Hay tres funciones públicas:
\texttt{get\_engine()} (singleton), \texttt{run\_sql\_file(path)}
(ejecuta un archivo en una sola transacción) y \texttt{read\_sql(sql)}
(devuelve un \texttt{pandas.DataFrame}). El listado completo aparece en
el apéndice~\ref{ap:python}.

\section{Logging y trazabilidad}

Todos los pasos de la ETL escriben a través del logger común
\texttt{src.utils.get\_logger}, que aplica un formato consistente:

\begin{lstlisting}[language=Python,
  caption={\texttt{src/utils/logger.py}: logger compartido.},
  label={lst:logger}]
def get_logger(name: str = "proyecto-final", level: int = logging.INFO) -> logging.Logger:
    logger = logging.getLogger(name)
    if logger.handlers:
        return logger
    logger.setLevel(level)
    h = logging.StreamHandler()
    h.setFormatter(logging.Formatter(
        "%(asctime)s [%(levelname)s] %(name)s — %(message)s",
        datefmt="%H:%M:%S"))
    logger.addHandler(h)
    logger.propagate = False
    return logger
\end{lstlisting}

Este logger se utiliza por \texttt{src.etl.load} y por el resto de
módulos analíticos para emitir trazas tipo
\texttt{HH:MM:SS [INFO] etl.load — Ejecutando 60\_load\_fact\_sales.sql}.

\section{Reproducibilidad}

El pipeline completo se reproduce desde cero con cuatro comandos. Todos
ellos se invocan sobre el entorno \texttt{UAX}:

\begin{lstlisting}[language=bash,
  caption={Reproducción end-to-end del pipeline.},
  label={lst:reproducir}]
# 1. Variables de entorno
cp .env.example .env   # rellenar DB_USER y DB_PASSWORD

# 2. Dependencias
conda run -n UAX pip install -r requirements.txt

# 3. Pipeline completo (ETL + analitica)
conda run -n UAX jupyter nbconvert --to notebook --execute --inplace notebooks/0*.ipynb

# 4. Dashboard
conda run -n UAX python -m dashboard.html.build
conda run -n UAX python -m dashboard.html.server --port 8001
\end{lstlisting}

\noindent El conjunto de tests de \texttt{pytest} actúa como red de
seguridad y se describe en el capítulo~\ref{ch:conclusiones}.

\chapter{CLTV y métricas de cliente}
\label{ch:cltv}

\section{Fórmula del CLTV}

El enunciado fija el CLTV (\emph{Customer Lifetime Value}) por cliente
mediante la siguiente expresión, válida sobre una ventana temporal $t$:

\begin{align}
  \mathrm{CLTV}(c)
  \;=\;
  \text{Ingresos}_t(c) \;\times\;
  \text{Margen}_t(c) \;\times\;
  \text{Frecuencia}_t(c) \;\times\;
  R_t(c).
  \label{eq:cltv}
\end{align}

En la implementación del proyecto, la ventana es la cobertura efectiva
del dataset, $\,2020\text{-}01 \to 2026\text{-}01$, esto es, $W=72$ meses
naturales.

\section{Operacionalización}

Cada componente de la ecuación~\eqref{eq:cltv} se traduce a una columna
calculada por cliente en
\texttt{data/processed/cltv.parquet}. La
tabla~\ref{tab:cltv-componentes} formaliza esa traducción.

\begin{table}[H]
\centering
\caption{Operacionalización de la fórmula de CLTV sobre el DWH.}
\label{tab:cltv-componentes}
\small
\begin{tabular}{l p{9.5cm}}
\toprule
\textbf{Componente} & \textbf{Definición} \\
\midrule
$\text{Ingresos}_t$  &
  $\displaystyle
    \sum_{l \in \text{ventas}(c)} \text{subtotal}_l
    \;-\;
    \sum_{r \in \text{devoluciones}(c)} \text{returned\_value}_r$
  \;(columna \texttt{net\_revenue}). \\
$\text{Margen}_t$    &
  $\dfrac{\text{Ingresos}_t - \sum_l \text{unit\_cost}_l \cdot
    \text{quantity}_l}{\text{Ingresos}_t}$
  \;(columna \texttt{margin\_rate}, proporción $\in [-\,\dots\,,1]$). \\
$\text{Frecuencia}_t$ &
  número de tickets distintos:
  $|\{\,\text{sale\_id} : c\,\}|$
  \;(columna \texttt{frequency}). \\
$R_t$                 &
  retención: meses con al menos una compra dividido por la ventana,
  $R_t = \texttt{active\_months}/W$ con $W=72$
  \;(columna \texttt{retention\_rate}). \\
\bottomrule
\end{tabular}
\end{table}

\noindent\textbf{Comentarios.} El margen es una \emph{proporción} sobre los
ingresos netos, no un valor monetario; multiplicado por $\text{Ingresos}_t$
recupera la magnitud monetaria. Los clientes con $\text{Ingresos}_t < 0$
(devoluciones que superan ventas) o con coste mayor que ingresos pueden
tener $\text{Margen}_t$ negativo o $> 1$, lo cual se preserva tal cual sin
truncar para no falsear el dato original (véase
capítulo~\ref{ch:conclusiones}, \emph{Limitaciones}).

\section{Implementación}

La operacionalización anterior se reparte en dos lugares: una consulta SQL
en el notebook~03 que agrega los componentes por cliente, y el módulo
\texttt{src/analytics/cltv.py} que aplica la fórmula multiplicativa de la
ecuación~\eqref{eq:cltv}. El módulo entero es deliberadamente compacto:

\begin{lstlisting}[language=Python,
  caption={\texttt{src/analytics/cltv.py}: aplicación de la fórmula de
  CLTV sobre las componentes ya agregadas.},
  label={lst:cltv_module}]
"""Calculo de CLTV por cliente.

Formula del enunciado:
    CLTV = Ingresos_t x Margen_t x Frecuencia_t x R_t

Donde, en la ventana temporal t:
    - Ingresos_t   = ingresos netos del cliente (ventas - devoluciones).
    - Margen_t     = (ingresos - coste) / ingresos  (proporcion).
    - Frecuencia_t = n. de tickets distintos del cliente.
    - R_t          = retencion = meses_activos / meses_de_la_ventana.
"""
from __future__ import annotations
import pandas as pd

def compute_cltv(df : pd.DataFrame) -> pd.DataFrame:
    """Calcula el CLTV de cada cliente con compras."""
    df["cltv"] = (
        df["net_revenue"].astype(float)
        * df["margin_rate"].astype(float)
        * df["frequency"].astype(float)
        * df["retention_rate"].astype(float)
    )
    return df.sort_values("cltv", ascending=False).reset_index(drop=True)
\end{lstlisting}

\noindent El test
\texttt{test\_cltv\_formula\_consistency} (apéndice~\ref{ap:tests})
recalcula la fórmula sobre las primeras $50$ filas y comprueba que el
valor almacenado coincide con el producto de las cuatro componentes.

\section{Resultados agregados}
\label{sec:cltv-resultados}

La tabla~\ref{tab:cltv-stats} resume las cinco distribuciones implicadas.
Los valores se han leído directamente de \texttt{cltv.parquet} con
\texttt{pandas.DataFrame.describe} y redondeado a unidades manejables.

\begin{table}[H]
\centering
\caption{Estadísticos por componente de CLTV ($n=5\,750$ clientes; ventana
$W=72$ meses).}
\label{tab:cltv-stats}
\small
\begin{tabular}{l r r r}
\toprule
\textbf{Magnitud} & \textbf{Mediana} & \textbf{Media} & \textbf{Máximo} \\
\midrule
Ingresos brutos (\texteuro)      & $149{,}90$  & $1\,683{,}26$  & $24\,922{,}00$ \\
Margen (proporción)             & $0{,}400$   & $0{,}364$      & $1{,}000$      \\
Frecuencia (tickets)            & $1$         & $3{,}48$       & $34$           \\
Retención $R_t$ ($\,/72$)       & $0{,}0139$  & $0{,}0436$     & $0{,}4028$     \\
\textbf{CLTV (\texteuro)}        & \textbf{0{,}72} & \textbf{$3\,165{,}48$} &
                                   \textbf{$109\,734{,}95$} \\
\bottomrule
\end{tabular}
\end{table}

\noindent La distancia entre mediana y media del CLTV ($0{,}72$\,\texteuro\
frente a $3\,165$\,\texteuro) refleja la asimetría extrema de la
distribución: la mayoría de clientes presenta frecuencia $1$ y retención
$\approx 1/72$, mientras que un pequeño grupo recurrente concentra todo
el valor (figura~\ref{fig:cltv-hist}). El máximo, $109\,735$\,\texteuro,
corresponde al cliente con $\text{customer\_id}=212$.

\begin{figure}[H]
  \centering
  \includegraphics[width=\linewidth]{cltv_distribution.png}
  \caption[Distribución del CLTV]{%
    Distribución del CLTV en \texttt{cltv.parquet}. Histograma en escala
    $\log_{10}$ a la izquierda; \emph{box-plot} con eje horizontal
    en escala simétrica logarítmica (\texttt{symlog}, $\text{linthresh}=1$)
    a la derecha.}
  \label{fig:cltv-hist}
\end{figure}

\section{Métricas adicionales por cliente}

Además del CLTV, el notebook~04 genera tres métricas complementarias en
\texttt{customer\_metrics.parquet}:

\begin{itemize}
  \item \textbf{Recencia} (\texttt{recency\_days}): días desde la última
        compra hasta la fecha de corte. Mediana $421$\,d, P75
        $\approx 1\,300$\,d, P95 $\approx 1\,782$\,d, máximo $1\,824$\,d
        (~$5$ años).
  \item \textbf{Ticket medio} o \emph{AOV} (\texttt{aov}): cociente entre
        ingresos netos y frecuencia. Mediana $129{,}99$\,\texteuro, media
        $210{,}04$\,\texteuro, P75 $\approx 300$\,\texteuro\ y máximo
        $1\,199{,}97$\,\texteuro\ (un AOV negativo aparece para los
        clientes con devolución superior a la venta).
  \item \textbf{Tasa de devolución} (\texttt{return\_rate}):
        $\texttt{returned\_value} / \texttt{gross\_revenue}$. Mediana $0$,
        media $0{,}069$, máximo $2{,}0$ (clientes en los que la cantidad
        devuelta duplica la cantidad vendida; este caso \emph{anómalo} es
        un artefacto del dataset sintético, véase
        capítulo~\ref{ch:conclusiones}).
\end{itemize}

\noindent La tabla~\ref{tab:metricas-extra} ofrece la lectura conjunta de
los estadísticos.

\begin{table}[H]
\centering
\caption{Estadísticos de las métricas adicionales por cliente
($n=5\,750$).}
\label{tab:metricas-extra}
\small
\begin{tabular}{l r r r r}
\toprule
\textbf{Métrica} & \textbf{Mediana} & \textbf{Media} & \textbf{P75} & \textbf{Máx.}\\
\midrule
Recencia (días)         & $421$    & $673{,}5$  & $1\,300$  & $1\,824$    \\
AOV (\texteuro)          & $129{,}99$ & $210{,}04$ & $299{,}99$ & $1\,199{,}97$ \\
Tasa de devolución     & $0$      & $0{,}069$  & $0$       & $2{,}0$     \\
\bottomrule
\end{tabular}
\end{table}

Los tests \texttt{test\_metrics\_aov\_matches\_revenue\_div\_frequency} y
\texttt{test\_recency\_is\_non\_negative} (apéndice~\ref{ap:tests})
formalizan las dos invariantes clave de estas métricas: el AOV recompone
los ingresos netos a partir de la frecuencia, y la recencia nunca puede
ser negativa.

\section{Segmentación: PCA + KMeans}
\label{sec:segmentacion}

La segmentación parte de un panel por cliente con ocho \emph{features}:
\texttt{net\_revenue}, \texttt{margin\_rate}, \texttt{frequency},
\texttt{retention\_rate}, \texttt{aov}, \texttt{recency\_days},
\texttt{return\_rate} y \texttt{cltv}. El pipeline es la composición
canónica $X\to Z=\text{StandardScaler}(X)\to Y=\text{PCA}_{n=2}(Z)\to
\text{KMeans}_{k=3}(Y)$, encapsulada en
\texttt{src/analytics/clustering.py}.

\subsection*{Fundamento matemático}

Sea $S=\tfrac{1}{n-1}Z^{\top}Z$ la covarianza muestral. Las direcciones
óptimas del PCA verifican

\begin{align}
  w_i \;=\; \argmax_{\|w\|=1}\; w^{\top}\,S\,w,
\end{align}

cuya solución por las condiciones KKT da $S\,w_i = \lambda_i\,w_i$: los
\emph{vectores propios} de $S$ ordenados por \emph{valor propio}
descendente. Las dos primeras componentes explican $57{,}1\%$ y
$23{,}2\%$ de la varianza ($\Sigma \approx 80{,}3\%$), valores leídos del
atributo \texttt{explained\_variance\_ratio\_} de
\texttt{models/pca.joblib}.

El KMeans minimiza la inercia intra-cluster

\begin{align}
  J(C_1,\dots,C_k)
  \;=\;
  \sum_{i=1}^{k}\;\sum_{y\in C_i}\;\bigl\|y - \mu_i\bigr\|_2^{2},
  \qquad
  \mu_i \;=\; \tfrac{1}{|C_i|}\sum_{y\in C_i}y,
  \label{eq:kmeans-inertia}
\end{align}

mediante el algoritmo de Lloyd, entrenado sobre el espacio reducido
$Y\in\mathbb{R}^{n\times 2}$. La curva del codo (figura~\ref{fig:elbow})
cae de $13\,360$ ($k=2$) a $2\,934$ ($k=3$) y desciende suavemente desde
ahí; se fija $k^{\star} = 3$.

\begin{figure}[H]
  \centering
  \includegraphics[width=0.65\linewidth]{elbow.png}
  \caption{Inercia $J^{\star}(k)$ frente a $k\in\{2,\dots,8\}$. El codo en
  $k^{\star}=3$ justifica la elección final.}
  \label{fig:elbow}
\end{figure}

\subsection*{Resultado: perfiles de cluster}

\begin{figure}[H]
  \centering
  \includegraphics[width=0.78\linewidth]{pca_clusters.png}
  \caption{Proyección PC1--PC2 de los $5\,750$ clientes coloreada por
  cluster KMeans.}
  \label{fig:pca-clusters}
\end{figure}

\begin{table}[H]
\centering
\caption{Perfil mediano por cluster, leído de
\texttt{models/customer\_segments.parquet}.}
\label{tab:cluster-perfil}
\renewcommand{\arraystretch}{1.25}
\resizebox{\linewidth}{!}{%
\begin{tabular}{l r r r r r r r r r}
\toprule
\textbf{Cluster} & \textbf{$n$} &
\texttt{net\_rev.} & \texttt{margin} & \texttt{freq.} &
\texttt{ret.} & \texttt{aov} & \texttt{recency} &
\texttt{ret\_rate} & \textbf{CLTV} \\
\midrule
0 (\emph{ocasional})   & $4\,576$ & $129{,}90$    & $0{,}400$ & $1$  & $0{,}014$ & $129{,}90$ & $854$ & $0{,}000$ & $0{,}72$        \\
1 (\emph{alto valor})  & $740$    & $11\,461{,}53$ & $0{,}388$ & $20$ & $0{,}250$ & $572{,}24$ & $82$  & $0{,}022$ & $21\,191{,}30$ \\
2 (\emph{devolutivo})  & $434$    & $0{,}00$       & $0{,}000$ & $1$  & $0{,}014$ & $0{,}00$   & $856$ & $1{,}000$ & $0{,}00$        \\
\bottomrule
\end{tabular}}
\end{table}

El \emph{cluster 0} ($\approx 80\%$ de la base) recoge los clientes
\emph{one-shot}. El \emph{cluster 1} ($\approx 13\%$) son recurrentes con
mediana de $20$ tickets, AOV de $572$\,\texteuro\ y retención del $25\%$,
y concentran el grueso del CLTV agregado. El \emph{cluster 2}
($\approx 7{,}5\%$) son clientes devolutivos con tasa de devolución
mediana del $100\%$ — el grupo separado por PC2, y ruido desde el punto
de vista de valor.

\chapter{Dashboard}
\label{ch:dashboard}

\section{Arquitectura}

El dashboard del proyecto es un \emph{sitio web estático} generado por
\texttt{dashboard/html/build.py} y servido por un pequeño backend Flask
(\texttt{dashboard/html/server.py}) que añade unos cuantos endpoints REST
sobre el DWH y los artefactos serializados. El build genera un único
fichero \texttt{index.html} (de aproximadamente $1{,}1$\,MB) con todas las
vistas embebidas y todo el JavaScript de interactividad inline; los datos
en vivo se obtienen vía \emph{fetch} contra los endpoints del servidor. El
servidor también detecta automáticamente si PostgreSQL está accesible y
conmuta entre dos modos:

\begin{itemize}
  \item \textbf{Modo \emph{vivo}}: las vistas de \emph{Visión General} y
        \emph{Productos} consultan directamente \texttt{dwh.fact\_sales}
        y \texttt{dwh.fact\_returns} para mostrar KPIs del periodo,
        evolución diaria y \emph{rankings} de tiendas y productos.
  \item \textbf{Modo \emph{offline}}: ante la ausencia de PostgreSQL,
        las vistas dependientes de DWH se reemplazan por su versión
        construida sólo desde los \emph{parquet} ya generados.
\end{itemize}

\section{Endpoints REST}

El backend expone los siguientes endpoints, todos en JSON:

\begin{lstlisting}[language=Python,
  caption={Endpoints del servidor del dashboard, en \texttt{server.py}.},
  label={lst:endpoints}]
GET /api/health                       # ping; n_customers, dwh online?
GET /api/customers/search             # busqueda + filtros sobre customer_segments
GET /api/customer/<customer_id>       # ficha completa: comparativas + radar + timeline
GET /api/cluster/<cluster_id>         # perfil estadistico + top clientes + top productos
GET /api/cluster/predict              # simulador: feature_vec -> cluster + (pc1,pc2)
GET /api/cluster/feature_ranges       # rangos (p01..p99) por feature para los sliders
\end{lstlisting}

\noindent
La búsqueda admite filtros por \texttt{cluster}, mínimo y máximo de CLTV,
mínimo y máximo de frecuencia, y orden por CLTV, recencia, frecuencia o
nombre. El simulador \texttt{/api/cluster/predict} carga los tres
artefactos \texttt{scaler/pca/kmeans.joblib} en memoria y aplica
exactamente la misma cadena utilizada en el entrenamiento, de manera que
las predicciones sobre clientes ficticios sean reproducibles.

\section{Las seis vistas}

La navegación del dashboard se ofrece como un sumario con seis capítulos,
en línea con la estética editorial del informe. Cada capítulo del
dashboard corresponde a un módulo bajo
\texttt{dashboard/html/views/}\,:

\begin{enumerate}
  \item \textbf{Visión General} (\texttt{overview.py}): KPIs del año,
        inventario verificado del DWH, evolución diaria con media móvil
        a $7$ días, mezcla por categoría, \emph{treemap} categoría/marca,
        estacionalidad y top 10 de tiendas
        (figura~\ref{fig:dash-overview}).
  \item \textbf{Clientes} (\texttt{clientes.py}): hallazgo principal
        (porcentaje de \emph{one-shot}), distribución del CLTV con
        histograma logarítmico, ECDF y densidad por cluster, curva de
        Lorenz con coeficiente de Gini, y top $20$ de clientes
        (figura~\ref{fig:dash-clientes}).
  \item \textbf{Segmentación} (\texttt{segmentacion.py}): tres tarjetas
        de cluster con KPI por cabeza, mapa PCA con \emph{hexbin},
        tamaño de cluster, varianza explicada, \emph{heatmap}
        $z$-score y simulador interactivo
        (figura~\ref{fig:dash-segmentacion}).
  \item \textbf{Productos} (\texttt{productos.py}): KPIs del periodo,
        top $15$ de productos por ingresos, descomposición de
        devoluciones por motivo, top productos por valor devuelto
        (figura~\ref{fig:dash-productos}).
  \item \textbf{Ficha Cliente} (\texttt{ficha.py}): buscador en vivo
        sobre $5\,750$ clientes con filtros por cluster, CLTV o
        frecuencia; al seleccionar, despliega cabecera, comparativas,
        radar de percentiles e historial de compras leído del DWH
        (figura~\ref{fig:dash-ficha}).
  \item \textbf{Metodología} (\texttt{metodologia.py}): mapa de
        KPI~$\to$~fórmula~$\to$~origen, listado del pipeline ETL, la
        propia justificación del $k=3$, varianza explicada y estadísticos
        clave (figura~\ref{fig:dash-metodologia}).
\end{enumerate}

\section{Capturas de pantalla}

Las capturas de las seis vistas se obtuvieron con el servidor Flask en
ejecución sobre \texttt{localhost:8001} y exportadas en alta resolución.

\begin{figure}[H]
  \centering
  \includegraphics[width=\linewidth]{dashboard_overview.png}
  \caption[Dashboard, Visión General]{Dashboard, vista
  \emph{Visión General}: KPIs del año en curso, inventario verificado del
  DWH y comienzo de la serie diaria de ingresos.}
  \label{fig:dash-overview}
\end{figure}

\begin{figure}[H]
  \centering
  \includegraphics[width=\linewidth]{dashboard_clientes.png}
  \caption[Dashboard, Clientes]{Dashboard, vista \emph{Clientes y CLTV}:
  hallazgo principal (un único \emph{ticket} en el $87\%$ de la base),
  histograma logarítmico del CLTV, curva de Lorenz con
  $\text{Gini}=0{,}907$ y comparación mediana--media por cluster.}
  \label{fig:dash-clientes}
\end{figure}

\begin{figure}[H]
  \centering
  \includegraphics[width=\linewidth]{dashboard_segmentacion.png}
  \caption[Dashboard, Segmentación]{Dashboard, vista
  \emph{Segmentación}: tarjetas-resumen por cluster, mapa PCA con
  contornos, justificación del $k=3$, varianza explicada y simulador
  interactivo del modelo.}
  \label{fig:dash-segmentacion}
\end{figure}

\begin{figure}[H]
  \centering
  \includegraphics[width=\linewidth]{dashboard_productos.png}
  \caption[Dashboard, Productos]{Dashboard, vista
  \emph{Productos y devoluciones}: KPIs del periodo, top $15$ por
  ingresos, distribución por motivo de devolución y mapa de hechos
  utilizados (\texttt{fact\_sales}, \texttt{fact\_returns},
  \texttt{dim\_product}, \texttt{dim\_return\_reason}).}
  \label{fig:dash-productos}
\end{figure}

\begin{figure}[H]
  \centering
  \includegraphics[width=\linewidth]{dashboard_ficha.png}
  \caption[Dashboard, Ficha Cliente]{Dashboard, vista
  \emph{Ficha Cliente}: buscador en vivo y panel detallado del cliente
  con radar de percentiles, tabla comparativa contra mediana del cluster
  y mediana global, e historial de compras agregado por categoría.}
  \label{fig:dash-ficha}
\end{figure}

\begin{figure}[H]
  \centering
  \includegraphics[width=\linewidth]{dashboard_metodologia.png}
  \caption[Dashboard, Metodología]{Dashboard, vista
  \emph{Metodología}: pipeline ETL, justificación del $k=3$ por curva del
  codo, varianza explicada y estadísticos clave del CLTV.}
  \label{fig:dash-metodologia}
\end{figure}

\section{Decisiones de diseño}

El dashboard se ha construido siguiendo una estética \emph{editorial} o
de \emph{data-storytelling} (próxima al estilo de los \emph{longreads} de
publicaciones como Pudding o Bloomberg) antes que un patrón clásico de
``rejilla de KPIs''. Las decisiones más relevantes:

\paragraph{Una columna principal con párrafos cortos.}
Cada gráfico se introduce con una pequeña descripción que orienta la
lectura, lo que evita que el lector tenga que reconstruir el contexto
visual por su cuenta.

\paragraph{Paletas de colores reservadas a la rentabilidad.}
El color de los gráficos no codifica categorías arbitrarias sino
métricas semánticamente comparables: el margen aparece en rampa
azul\,$\to$\,verde, la concentración (Lorenz) en azul, y los clusters en
una terna estable (azul, ámbar, rojo) que se repite en todas las vistas
de segmentación.

\paragraph{Modo \emph{vivo} vs. \emph{offline} con misma maqueta.}
El sitio detecta la disponibilidad del DWH al arrancar y conmuta los
componentes pertinentes sin alterar la maqueta, lo que permite usarlo
como prototipo de presentación incluso cuando no hay base de datos
disponible.

\paragraph{Plotly inline + Flask para el detalle.}
Los gráficos se generan en el servidor mediante \texttt{plotly.io.to\_html}
y se embeben directamente en \texttt{index.html}, lo que evita
dependencias en tiempo de carga; los componentes verdaderamente
interactivos (\emph{ficha de cliente}, \emph{simulador}, \emph{cluster
detail}) son los que consumen los endpoints REST.

\section{Conclusiones}
\label{sec:conclusiones}

\subsection*{Hallazgos cuantitativos}

Sobre $5\,750$ clientes, $20\,000$ tickets y $42\,555$ líneas de venta, el
pipeline ETL puebla un DWH en estrella ($6$~dimensiones, $2$~hechos) de
forma reproducible: una sola invocación de
\texttt{jupyter nbconvert --execute notebooks/0*} regenera los esquemas
\texttt{staging} y \texttt{dwh}, los \emph{parquet} de analítica y los
modelos de \texttt{models/}. La distribución del CLTV es extremadamente
asimétrica (mediana $0{,}72$\,\texteuro\ frente a media $3\,165$\,\texteuro,
máximo $109\,735$\,\texteuro), con un coeficiente de Gini~$\approx 0{,}91$:
el $87\%$ de los clientes son \emph{one-shot} y el cluster ``alto valor''
($13\%$ de la base) concentra el grueso del CLTV agregado. El segmentador
PCA(2) $+$ KMeans(3) recupera tres clusters interpretables (ocasionales,
alto valor, devolutivos) con $80{,}3\%$ de varianza explicada y un codo
claro en $k^{\star}=3$.

\subsection*{Limitaciones}

El dataset es \emph{sintético} y algunas decisiones de diseño preservan los
datos tal cual aunque generen artefactos: existen $10$ clientes con
$\texttt{return\_rate} > 1$ (devoluciones que superan la venta original) y
$1$ producto sin \texttt{unit\_cost} en \texttt{central\_product}; ambos
casos se conservan y se gestionan con \texttt{coalesce(\dots, 0)} en el
cálculo de coste. La retención $R_t$ es una proxy discreta
($\texttt{active\_months}/72$) y el KMeans se entrena sobre el espacio
reducido por PCA(2), que descarta un $\sim 20\%$ de la varianza original.

\subsection*{Líneas futuras y validación}

Las extensiones naturales son: ampliar las \emph{features} con señales
territoriales (\texttt{district}, \texttt{area\_type}) y promocionales
(\texttt{discount\_percent}); validar $k$ con métricas de silueta y
Davies--Bouldin; sustituir la fórmula multiplicativa por modelos
probabilísticos tipo BG/NBD + Gamma--Gamma; y refactorizar el pipeline
para ingestión incremental con \emph{slowly changing dimensions}. La
calidad queda anclada por nueve pruebas \texttt{pytest} en
\texttt{tests/test\_cltv.py} y \texttt{tests/test\_dwh\_integrity.py}
(integridad referencial, unicidad de \emph{business keys}, conservación de
ingresos \texttt{staging} $\to$ \texttt{dwh}, y consistencia algebraica de
la fórmula de CLTV).


% --------------------------------------------------------------- Apéndice ---
\appendix
\chapter{Apéndice}
\label{ch:apendice}

\section{Esquema completo del DWH}
\label{ap:ddl}

Se reproducen aquí los \texttt{DDL} de \texttt{sql/02\_dwh\_schema/} tal
como están en el repositorio. Las claves subrogadas usan \texttt{serial}
y las claves de negocio se marcan \texttt{unique}.

\begin{lstlisting}[language=SQLpostgres,
  caption={\texttt{sql/02\_dwh\_schema/00\_create\_schema.sql}: creación
  del esquema \texttt{dwh}.},
  label={lst:ap-create-schema}]
-- Esquema dwh: modelo en estrella.
create schema if not exists dwh;
\end{lstlisting}

\begin{lstlisting}[language=SQLpostgres,
  caption={\texttt{sql/02\_dwh\_schema/10\_dim\_tables.sql}: dimensiones
  del DWH.},
  label={lst:ap-dim-tables}]
-- Tablas de dimension.

drop table if exists dwh.dim_date cascade;
create table dwh.dim_date (
    date_id      int           primary key,         -- yyyymmdd
    full_date    date          not null unique,
    year         smallint      not null,
    quarter      smallint      not null,
    month        smallint      not null,
    month_name   varchar(15)   not null,
    day          smallint      not null,
    day_of_week  smallint      not null,
    day_name     varchar(15)   not null,
    week_of_year smallint      not null,
    is_weekend   boolean       not null
);

drop table if exists dwh.dim_customer cascade;
create table dwh.dim_customer (
    customer_sk  serial        primary key,
    customer_id  int           not null unique,
    full_name    varchar(200),
    email        varchar(200),
    phone        varchar(50),
    signup_date  date
);

drop table if exists dwh.dim_product cascade;
create table dwh.dim_product (
    product_sk   serial        primary key,
    product_id   int           not null unique,
    name         varchar(200),
    category     varchar(100),
    brand        varchar(100),
    sku          varchar(50),
    unit_cost    numeric(12,2),
    unit_price   numeric(12,2)
);

drop table if exists dwh.dim_store cascade;
create table dwh.dim_store (
    store_sk     serial        primary key,
    store_id     int           not null unique,
    name         varchar(200),
    address      varchar(200),
    city         varchar(100),
    postal_code  varchar(10),
    district     varchar(100),
    area_type    varchar(50),
    zone_orientation varchar(50),
    opened_date  date
);

drop table if exists dwh.dim_offer cascade;
create table dwh.dim_offer (
    offer_sk         serial      primary key,
    offer_id         int         not null unique,
    name             varchar(200),
    discount_percent numeric(5,2),
    start_date       date,
    end_date         date
);

drop table if exists dwh.dim_return_reason cascade;
create table dwh.dim_return_reason (
    reason_sk    serial        primary key,
    reason_id    int           not null unique,
    reason       text,
    active       boolean
);
\end{lstlisting}

\begin{lstlisting}[language=SQLpostgres,
  caption={\texttt{sql/02\_dwh\_schema/20\_fact\_tables.sql}: hechos del
  DWH.},
  label={lst:ap-fact-tables}]
-- Tablas de hechos.

drop table if exists dwh.fact_sales cascade;
create table dwh.fact_sales (
    sale_item_id  int           primary key,
    sale_id       int           not null,
    sale_date_id  int           not null references dwh.dim_date(date_id),
    customer_sk   int           references dwh.dim_customer(customer_sk),
    product_sk    int           references dwh.dim_product(product_sk),
    store_sk      int           references dwh.dim_store(store_sk),
    offer_sk      int           references dwh.dim_offer(offer_sk),
    quantity      int           not null,
    unit_price    numeric(12,2) not null,
    unit_cost     numeric(12,2),
    subtotal      numeric(14,2) not null,
    cost          numeric(14,2),
    margin        numeric(14,2)
);
create index ix_fs_customer on dwh.fact_sales(customer_sk);
create index ix_fs_product  on dwh.fact_sales(product_sk);
create index ix_fs_store    on dwh.fact_sales(store_sk);
create index ix_fs_date     on dwh.fact_sales(sale_date_id);

drop table if exists dwh.fact_returns cascade;
create table dwh.fact_returns (
    return_id        int           primary key,
    sale_item_id     int           references dwh.fact_sales(sale_item_id),
    return_date_id   int           not null references dwh.dim_date(date_id),
    customer_sk      int           references dwh.dim_customer(customer_sk),
    product_sk       int           references dwh.dim_product(product_sk),
    store_sk         int           references dwh.dim_store(store_sk),
    reason_sk        int           references dwh.dim_return_reason(reason_sk),
    quantity         int           not null,
    returned_value   numeric(14,2)
);
create index ix_fr_customer on dwh.fact_returns(customer_sk);
create index ix_fr_product  on dwh.fact_returns(product_sk);
create index ix_fr_date     on dwh.fact_returns(return_date_id);
\end{lstlisting}

\section{Transformaciones SQL (resumen)}
\label{ap:transformaciones}

Las transformaciones \texttt{20\_load\_dim\_customer.sql} a
\texttt{50\_load\_dim\_misc.sql} se reproducen tal cual; las dos cargas
de hechos ya aparecen en los listados~\ref{lst:fact_sales_load} y los
\texttt{LEFT JOIN} de \texttt{fact\_returns} se citan a continuación.

\begin{lstlisting}[language=SQLpostgres,
  caption={\texttt{30\_load\_dim\_product.sql}: combinación de
  \texttt{product} y \texttt{central\_product}.},
  label={lst:ap-dim-product}]
-- dim_product: combina staging.product (50) con staging.central_product (49).
truncate table dwh.dim_product restart identity cascade;

insert into dwh.dim_product (product_id, name, category, brand, sku, unit_cost, unit_price)
select  p.product_id,
        coalesce(cp.name, p.name)                                         as name,
        coalesce(cat.name, p.category)                                    as category,
        coalesce(b.name,  p.manufacturer)                                 as brand,
        cp.sku,
        cp.unit_cost,
        coalesce(cp.unit_price, p.price)                                  as unit_price
from staging.product p
left join staging.central_product cp on cp.product_id = p.product_id
left join staging.category   cat on cat.category_id = cp.category_id
left join staging.brand      b   on b.brand_id      = cp.brand_id;
\end{lstlisting}

\begin{lstlisting}[language=SQLpostgres,
  caption={\texttt{70\_load\_fact\_returns.sql}: reutilización de FKs de
  \texttt{fact\_sales}.},
  label={lst:ap-fact-returns}]
truncate table dwh.fact_returns cascade;

insert into dwh.fact_returns
       (return_id, sale_item_id, return_date_id, customer_sk, product_sk,
        store_sk, reason_sk, quantity, returned_value)
select  ri.return_id,
        ri.sale_item_id,
        to_char(ri.return_date, 'YYYYMMDD')::int  as return_date_id,
        fs.customer_sk,
        fs.product_sk,
        fs.store_sk,
        drr.reason_sk,
        ri.quantity,
        ri.quantity * fs.unit_price               as returned_value
from staging.return_item ri
left join dwh.fact_sales       fs  on fs.sale_item_id = ri.sale_item_id
left join dwh.dim_return_reason drr on drr.reason_id  = ri.reason_id;
\end{lstlisting}

\section{Listados completos de Python}
\label{ap:python}

\begin{lstlisting}[language=Python,
  caption={\texttt{src/analytics/cltv.py}.},
  label={lst:ap-cltv}]
"""Calculo de CLTV por cliente.

Formula del enunciado:
    CLTV = Ingresos_t x Margen_t x Frecuencia_t x R_t

Donde, en la ventana temporal t:
    - Ingresos_t   = ingresos netos del cliente (ventas - devoluciones).
    - Margen_t     = (ingresos - coste) / ingresos  (proporcion).
    - Frecuencia_t = n. de tickets distintos del cliente.
    - R_t          = retencion = meses_activos / meses_de_la_ventana.
"""
from __future__ import annotations
import pandas as pd

def compute_cltv(df : pd.DataFrame) -> pd.DataFrame:
    """Calcula el CLTV de cada cliente con compras."""
    df["cltv"] = (
        df["net_revenue"].astype(float)
        * df["margin_rate"].astype(float)
        * df["frequency"].astype(float)
        * df["retention_rate"].astype(float)
    )
    return df.sort_values("cltv", ascending=False).reset_index(drop=True)
\end{lstlisting}

\begin{lstlisting}[language=Python,
  caption={\texttt{src/analytics/clustering.py}.},
  label={lst:ap-clustering}]
"""Segmentacion de clientes con PCA + KMeans."""
from __future__ import annotations

from dataclasses import dataclass
from pathlib import Path

import numpy as np
import pandas as pd
from sklearn.cluster import KMeans
from sklearn.decomposition import PCA
from sklearn.preprocessing import StandardScaler


@dataclass
class ClusteringResult:
    df: pd.DataFrame                  # df + columnas pc1, pc2, cluster
    pca: PCA
    scaler: StandardScaler
    kmeans: KMeans
    explained_variance: np.ndarray
    inertia_curve: dict[int, float]   # k -> inertia, para escoger k


def _elbow(X: np.ndarray, ks=range(2, 9), random_state=42) -> dict[int, float]:
    return {k: KMeans(n_clusters=k, n_init=10, random_state=random_state).fit(X).inertia_
            for k in ks}


def run_pca_kmeans(
    df: pd.DataFrame,
    feature_cols: list[str],
    n_clusters: int = 3,
    n_components: int = 2,
    random_state: int = 42,
) -> ClusteringResult:
    """Estandariza, aplica PCA(n_components) y KMeans(n_clusters)."""
    X = df[feature_cols].astype(float).fillna(0).to_numpy()
    scaler = StandardScaler().fit(X)
    Xs = scaler.transform(X)

    pca = PCA(n_components=n_components, random_state=random_state).fit(Xs)
    Xp = pca.transform(Xs)

    inertia = _elbow(Xp, random_state=random_state)
    km = KMeans(n_clusters=n_clusters, n_init=10, random_state=random_state).fit(Xp)

    out = df.copy()
    for i in range(n_components):
        out[f"pc{i+1}"] = Xp[:, i]
    out["cluster"] = km.labels_
    return ClusteringResult(out, pca, scaler, km, pca.explained_variance_ratio_, inertia)


def save_artifacts(result: ClusteringResult, out_dir: str | Path) -> None:
    import joblib
    out = Path(out_dir)
    out.mkdir(parents=True, exist_ok=True)
    joblib.dump(result.scaler, out / "scaler.joblib")
    joblib.dump(result.pca, out / "pca.joblib")
    joblib.dump(result.kmeans, out / "kmeans.joblib")
    result.df.to_parquet(out / "customer_segments.parquet", index=False)
\end{lstlisting}

\begin{lstlisting}[language=Python,
  caption={\texttt{src/db/engine.py}.},
  label={lst:ap-engine}]
"""Conexion a Postgres y helpers para ejecutar SQL."""
from __future__ import annotations

from functools import lru_cache
from pathlib import Path

import pandas as pd
from sqlalchemy import Engine, create_engine, text

from config.settings import settings


@lru_cache(maxsize=1)
def get_engine() -> Engine:
    """Engine SQLAlchemy reutilizable (singleton)."""
    return create_engine(settings.database_url, future=True, pool_pre_ping=True)


def run_sql(sql: str, **params) -> None:
    """Ejecuta una sentencia SQL (DDL o DML) con autocommit."""
    eng = get_engine()
    with eng.begin() as conn:
        conn.execute(text(sql), params or None)


def run_sql_file(path: str | Path) -> None:
    """Ejecuta un archivo .sql completo. Asume sentencias separadas por ';'.

    Compatible con bloques DO $$ ... $$ porque ejecuta el archivo entero
    en una sola transaccion usando exec_driver_sql.
    """
    raw = Path(path).read_text(encoding="utf-8")
    eng = get_engine()
    with eng.begin() as conn:
        conn.exec_driver_sql(raw)


def read_sql(sql: str, **params) -> pd.DataFrame:
    """Devuelve un DataFrame desde una consulta."""
    eng = get_engine()
    return pd.read_sql(text(sql), eng, params=params or None)
\end{lstlisting}

\begin{lstlisting}[language=Python,
  caption={\texttt{src/etl/load.py}.},
  label={lst:ap-load}]
"""Helpers de carga: ejecutar SQL en orden y resetear esquemas."""
from __future__ import annotations

from pathlib import Path

from src.db import run_sql, run_sql_file
from src.utils import get_logger

log = get_logger("etl.load")


def run_sql_directory(directory: str | Path) -> list[Path]:
    """Ejecuta los .sql de un directorio en orden alfabetico."""
    files = sorted(Path(directory).glob("*.sql"))
    for f in files:
        log.info("Ejecutando %s", f.name)
        run_sql_file(f)
    return files
\end{lstlisting}

\section{Manifiesto de tests}
\label{ap:tests}

Los listados completos de \texttt{tests/test\_cltv.py} y
\texttt{tests/test\_dwh\_integrity.py} aparecen a continuación. En el
capítulo~\ref{ch:conclusiones} ya se describe brevemente cada test.

\begin{lstlisting}[language=Python,
  caption={\texttt{tests/test\_cltv.py}.},
  label={lst:ap-test-cltv}]
"""Validan la logica de CLTV y metricas."""
import pytest

from src.db import read_sql

try:
    read_sql("select 1")
except Exception as e:  # pragma: no cover
    pytest.skip(f"DB no disponible: {e}", allow_module_level=True)

from src.analytics import compute_cltv, compute_customer_metrics


def test_cltv_columns_present():
    df = compute_cltv()
    expected = {
        "customer_id", "full_name", "gross_revenue", "net_revenue",
        "margin_rate", "frequency", "retention_rate", "cltv",
    }
    assert expected.issubset(df.columns)


def test_cltv_formula_consistency():
    df = compute_cltv().head(50).copy()
    expected = (
        df["net_revenue"].astype(float)
        * df["margin_rate"].astype(float)
        * df["frequency"].astype(float)
        * df["retention_rate"].astype(float)
    )
    assert ((df["cltv"] - expected).abs() < 1e-6).all()


def test_metrics_aov_matches_revenue_div_frequency():
    m = compute_customer_metrics().head(100).copy()
    pos = m[m["frequency"] > 0]
    expected = pos["net_revenue"].astype(float) / pos["frequency"].astype(float)
    assert ((pos["aov"].astype(float) - expected).abs() < 1e-6).all()


def test_recency_is_non_negative():
    m = compute_customer_metrics()
    assert (m["recency_days"] >= 0).all()
\end{lstlisting}

\begin{lstlisting}[language=Python,
  caption={\texttt{tests/test\_dwh\_integrity.py}.},
  label={lst:ap-test-dwh}]
"""Validan que el DWH este coherente respecto al raw."""
import pytest

from src.db import read_sql

# Si no hay BD disponible, saltamos todos los tests de integridad.
try:
    read_sql("select 1 as ok")
except Exception as e:  # pragma: no cover
    pytest.skip(f"DB no disponible: {e}", allow_module_level=True)


def test_fact_sales_matches_raw_sale_item_count():
    raw = read_sql("select count(*) as n from public.sale_item").iloc[0]["n"]
    dwh = read_sql("select count(*) as n from dwh.fact_sales").iloc[0]["n"]
    assert int(raw) == int(dwh)


def test_fact_returns_matches_raw_return_item_count():
    raw = read_sql("select count(*) as n from public.return_item").iloc[0]["n"]
    dwh = read_sql("select count(*) as n from dwh.fact_returns").iloc[0]["n"]
    assert int(raw) == int(dwh)


def test_dim_customer_unique_business_key():
    df = read_sql(
        "select customer_id, count(*) c from dwh.dim_customer group by 1 having count(*) > 1"
    )
    assert df.empty


def test_no_dangling_dim_date_in_facts():
    bad = read_sql("""
        select count(*) as n from dwh.fact_sales f
        where not exists (select 1 from dwh.dim_date d where d.date_id = f.sale_date_id)
    """).iloc[0]["n"]
    assert int(bad) == 0


def test_revenue_matches_between_staging_and_dwh():
    stg = read_sql("select sum(subtotal) as v from staging.sale_item").iloc[0]["v"]
    dwh = read_sql("select sum(subtotal) as v from dwh.fact_sales").iloc[0]["v"]
    assert abs(float(stg) - float(dwh)) < 0.01
\end{lstlisting}


\end{document}
